
\subsection{纳入算子 $\mathfrak q$}
\label{subsec:5-4-5}

本小节把典范模型的讨论与滤过
$A_\infty$ 代数
$(H(L;\Lambda_{0,\mathrm{nov}}),\mathfrak m)$
的 bulk 形变结合起来。
这里取 $\mathbb Q$ 系数。
令
$b=\operatorname{rank}H^{\mathrm{even}}(M;\mathbb Q)$，
选取一组基 $c_i$，
并令 $2d_i=\deg c_i$。
按次数配平，记
$\mathbf c_i=e^{1-d_i}c_i
\in H^2(M;\Lambda_{0,\mathrm{nov}})$。
\translatornote{原书实体页 356
写作 $e^{d_i-1}c_i\in
H^2(M;\Lambda^+_{0,\mathrm{nov}})$。
按约定 Conv.4 的 $\deg e=2$，
前者除 $d_i=1$ 外不为二次；
即使修正指数，所得类仍无正 $T$ 能量，
不能仅凭本处定义断言其属于 $\Lambda^+$。
见疑误 A-0339。}
选取形式参数 $r_1,\ldots,r_b$，并置
\[
\mathcal R=\mathbb Q[[r_1,\ldots,r_b]],
\qquad
\Lambda_{0,\mathrm{nov}}(\mathcal R)
=\Lambda_{0,\mathrm{nov}}
\otimes_{\mathbb Q}\mathcal R.
\]
将用算子 $\mathfrak q$
在 $\Lambda_{0,\mathrm{nov}}(\mathcal R)$ 上
定义滤过 $A_\infty$ 代数
$(C(L;\Lambda_{0,\mathrm{nov}}(\mathcal R)),
\mathfrak m^{\mathcal R})$。\liftIndex{I215}

令
$\mathbf x\in
B_kC(L;\Lambda_{0,\mathrm{nov}}(\mathcal R))[1]$。
\begin{definition}{5.4.57}\label{def:5-4-57}
定义
\begin{equation}
\mathfrak m_k^{\mathcal R}(\mathbf x)
=
\mathfrak q\left(
\exp\left(\sum_{i=1}^{b}r_i\mathbf c_i\right);
\mathbf x
\right).
\tag{5.4.58}\label{eq:5-4-58}
\end{equation}
\end{definition}
式 (5.4.58) 右端还可写为
\[
\sum_{a\ge0}
\sum_{i_1,\ldots,i_a=1}^{b}
r_{i_1}\cdots r_{i_a}\,
\mathfrak q
(\mathbf c_{i_1}\cdots\mathbf c_{i_a};\mathbf x)
\in C(L;\Lambda_{0,\mathrm{nov}}(\mathcal R)).
\]
以显然方式把
$\mathfrak m_k^{\mathcal R}$ 扩张为
\[
\mathfrak m_k^{\mathcal R}\colon
B_kC(L;\Lambda_{0,\mathrm{nov}}(\mathcal R))[1]
\longrightarrow
C(L;\Lambda_{0,\mathrm{nov}}(\mathcal R)).
\]

\begin{lemma}{5.4.59}\label{lem:5-4-59}
$(C(L;\Lambda_{0,\mathrm{nov}}(\mathcal R)),
\mathfrak m^{\mathcal R})$
是一个滤过 $A_\infty$ 代数。
\end{lemma}
其证明与引理 3.8.39 的证明相同。

\begin{proposition}{5.4.60}\label{prop:5-4-60}
滤过 $A_\infty$ 代数
$(C(L;\Lambda_{0,\mathrm{nov}}(\mathcal R)),
\mathfrak m^{\mathcal R})$
在同伦等价意义下不依赖于构造中的各项选择。
若 $\psi\colon M\to M'$ 是辛微分同胚，
则它诱导同伦等价
\[
\psi_*\colon
(C(L;\Lambda_{0,\mathrm{nov}}(\mathcal R)),
\mathfrak m^{\mathcal R})
\longrightarrow
(C(\psi(L);\Lambda_{0,\mathrm{nov}}(\mathcal R)),
\mathfrak m^{\mathcal R}).
\]
应注意，源侧的 $\mathcal R$ 由
$H^{\mathrm{even}}(M)$ 定义，
靶侧则由 $H^{\mathrm{even}}(M')$ 定义；
二者经 $\psi^{-1*}$ 识别。
\end{proposition}
证明只需略微修改定理 4.6.52 的证明。

现在考虑
$(C(L;\Lambda_{0,\mathrm{nov}}(\mathcal R)),
\mathfrak m^{\mathcal R})$
的典范模型。
置
\[
\Lambda^+_{0,\mathrm{nov}}(\mathcal R)
=\Lambda^+_{0,\mathrm{nov}}\otimes_{\mathbb Q}\mathcal R.
\]
这是
$\Lambda_{0,\mathrm{nov}}(\mathcal R)$
的一个理想，但不是极大理想；并且
\[
\Lambda_{0,\mathrm{nov}}(\mathcal R)
/\Lambda^+_{0,\mathrm{nov}}(\mathcal R)
\cong\mathcal R[e,e^{-1}].
\]

\begin{lemma}{5.4.61}\label{lem:5-4-61}
设
$(C(L;\mathcal R),
\overline{\mathfrak m}^{\mathcal R})$
为从
$(C(L;\Lambda_{0,\mathrm{nov}}(\mathcal R)),
\mathfrak m^{\mathcal R})$
诱导的 $\mathcal R$ 上的 $A_\infty$ 代数。
则有 $A_\infty$ 代数同构
\[
(C(L;\mathcal R),
\overline{\mathfrak m}^{\mathcal R})
\cong
(\overline C(L;\mathbb Q),\overline{\mathfrak m})
\otimes_{\mathbb Q}\mathcal R[e,e^{-1}].
\]
\end{lemma}
\begin{proof}
原书认为这直接来自同余
\translatornote{这一步的充分条件尚未由原书给出：
式 (3.8.34.1) 允许
$\mathfrak q_{1,0}(\mathbf c_i)$
有零 $T$ 能量分量，
而 $r_i$ 在本处并未获赋正 $T$ 能量。
因此原书的同余及本引理的这一证明
在目前假设下一般不能直接推出；
译稿保留结论与源书证明思路，
不臆补额外消失假设。
见疑误 A-0339。}
\[
\mathfrak q\left(
\exp\left(\sum_i r_i\mathbf c_i\right);
\mathbf x
\right)
\equiv\mathfrak m(\mathbf x)
\pmod{\Lambda^+_{0,\mathrm{nov}}(\mathcal R)}.
\]
\end{proof}

\begin{corollary}{5.4.62}\label{cor:5-4-62}
引理 5.4.28 对
$(C(L;\mathcal R),
\overline{\mathfrak m}^{\mathcal R})$
也成立。
\end{corollary}
在引理 5.4.28 的证明中，
曾用到 $R$ 是域；
推论 5.4.62 说，
即使这里的 $\mathcal R$ 不是域，
所需结论仍成立。
如注 5.4.29 所说，
构造典范模型时只有这一处用到
$R$ 为域的假设。
因此得到如下结果。

\begin{proposition}{5.4.63}\label{prop:5-4-63}
滤过 $A_\infty$ 代数
$(C(L;\Lambda_{0,\mathrm{nov}}(\mathcal R)),
\mathfrak m^{\mathcal R})$
存在含幺的典范模型；
其同构类不依赖于构造中的各项选择。
\end{proposition}

命题 5.4.63 给出一个超仓西映射。
沿用式 (5.4.9) 的
$\mathbf v_j,\mathbf w_j,\pi_i,s_j,h_j$
及其分量编号，
令
\[
z=\sum_{j=1}^{m}s_j\mathbf v_j
+\sum_{j=1}^{\ell}h_j\mathbf w_j.
\]
对 $1\le i\le m+\ell$，
有\translatornote{原书实体页 357
在第二个指数中把
$\sum_jh_j\mathbf w_j$
写成 $\sum_jh_j\mathbf w_i$；
这里按同式第一行及 (5.4.9)
恢复求和下标。
见疑误 A-0337。}
\[
\begin{aligned}
P_i&\in
\Lambda_{0,\mathrm{nov}}
[[r_1,\ldots,r_b;
s_1,\ldots,s_m;
h_1,\ldots,h_\ell]],\\
P_i
&=(\pi_i\circ\mathfrak m^{\mathcal R})(\exp z)\\
&=(\pi_i\circ\mathfrak q)
\left(
\exp\left(\sum_{a=1}^{b}r_a\mathbf c_a\right);
\exp z
\right).
\end{aligned}
\]
令
$P_i^{\mathrm{bos}}=
P_i|_{h_1=\cdots=h_\ell=0}$。
定义\translatornote{原书实体页 357
仅列 $i=1,\ldots,\ell$，
又以 $(P_1,\ldots,P_m)$ 作纯偶商；
但 (5.4.9) 规定前 $m$ 项是奇基分量，
定义 5.4.11 的纯偶理想
取后 $\ell$ 项。
此处依既定分量编号写全
$1\le i\le m+\ell$，
并以
$(P_{m+1}^{\mathrm{bos}},
\ldots,P_{m+\ell}^{\mathrm{bos}})$
作商。见疑误 A-0338。}
\[
\mathfrak M_{\mathrm{def}}
(C_{\mathrm{can}}(L;\mathcal R),
\mathfrak m^{\mathcal R})^{\mathrm{bos}}
=
\frac{
\Lambda_{0,\mathrm{nov}}
[[r_1,\ldots,r_b;s_1,\ldots,s_m]]
}{
(P_{m+1}^{\mathrm{bos}},
\ldots,P_{m+\ell}^{\mathrm{bos}})
}.
\]
它是交换 $\mathcal R$-代数，
且其作为 $\mathcal R$-代数的同构类
是这个拉格朗日子流形的不变量。

下述结论可用命题 5.4.15
的同样方法证明。
\begin{lemma}{5.4.64}\label{lem:5-4-64}
\[
\operatorname{Hom}\left(
\mathfrak M_{\mathrm{def}}
(C_{\mathrm{can}}(L;\mathcal R),
\mathfrak m^{\mathcal R})^{\mathrm{bos}},
\Lambda_{0,\mathrm{nov}}
\right)
\cong
\widehat{\mathcal M}(L)_{\mathrm{def}}.
\]
\end{lemma}

由这一纯偶代数所得的形式概形
\[
\operatorname{Spf}\left(
\mathfrak M_{\mathrm{def}}
(C_{\mathrm{can}}(L;\mathcal R),
\mathfrak m^{\mathcal R})^{\mathrm{bos}}
\right)
\]
是 $\mathcal R$ 上的形式概形。
\translatornote{原书实体页 358
在此 $\operatorname{Spf}$ 的自变量中
省去 $\mathrm{bos}$ 上标及结构记号；
这里依紧前的交换代数和
引理 5.4.64 补明。
见疑误 A-0340。}
这与同调镜像对称相容。
也可用同样方法证明“def, weak”版本。

还可以使用算子 $\mathfrak r$
来研究有两个拉格朗日子流形的情形：
可以在
$\Lambda_{0,\mathrm{nov}}(\mathcal R)$
上构造滤过 $A_\infty$ 双模的典范模型。
细节留待读者完成。
